\chapter{系统总体设计}

在明确研究问题与相关技术基础之后，本文需要进一步回答方法如何在系统层面被组织、各模块如何协同工作以及实验平台如何支持双主线对比的问题。因此，本章从整体架构角度出发，对工程图纸检索系统的目标、层次结构、模块划分与运行流程进行说明，为后续方法实现与实验复现奠定系统基础。

\section{系统目标}

本文系统的目标是在输入一张工程图纸后，从图库中检索出与其结构和语义相似的其他图纸，并返回排序分数、结构分析信息、多模态辅助信号和可解释性可视化结果。系统既要满足算法研究的可验证性，也要满足工程部署的可运行性。

\section{系统功能需求与性能需求}

结合任务书和开题报告中对毕业设计的功能描述，本文系统的需求可以划分为功能性需求与性能性需求两个方面。

\subsection{功能性需求}

\begin{enumerate}
  \item \textbf{图纸上传与检索需求。} 用户应能够上传查询图纸，并获得与之结构或语义相近的候选图纸结果列表。
  \item \textbf{图库管理需求。} 系统应支持图库统计、类别浏览、图像入库与索引重建等基础管理功能。
  \item \textbf{结果解释需求。} 除返回相似度分数外，系统还应提供结构相似信息、文本辅助信息以及可解释性分析结果。
  \item \textbf{实验支撑需求。} 系统不仅服务于在线检索，还应支持离线评测、模式切换和实验结果复查。
\end{enumerate}

\subsection{性能性需求}

\begin{enumerate}
  \item \textbf{检索精度要求。} 系统需要在多类别工程图纸场景下保持较高的前列命中率与排序稳定性，以满足工程复用需求。
  \item \textbf{检索效率要求。} 在单次查询过程中，系统应尽量压缩图像加载、特征提取、向量召回与结果组织的总体耗时。
  \item \textbf{可扩展性要求。} 系统应支持不同检索模式、不同图库规模以及后续新模块的接入。
  \item \textbf{可复现性要求。} 系统应能够通过统一脚本完成建库、评测与结果导出，保证实验过程可追踪、可复查。
\end{enumerate}

\section{系统总体架构}

系统整体由图纸输入层、路线选择层、特征提取层、向量检索层、重排层和结果分析层组成。考虑到本文已经将方法重构为两条主线，系统在进入正式检索前会先根据当前模式决定走“YOLO 清洗路线”还是“无清洗注意力路线”。整体处理流程如下：

\begin{enumerate}
  \item 对输入图纸进行格式统一，并按当前模式决定是否执行 YOLO 清洗与保留掩码生成。
  \item 使用 CLIP 视觉编码器提取图像嵌入，并根据路线不同执行基础视觉编码或掩码引导的局部聚合增强。
  \item 从图纸中提取结构描述子，并在需要时提取 OCR 文本描述子。
  \item 将视觉向量送入 Qdrant 完成候选召回。
  \item 对候选结果进行结构重排，并在需要时于前排阶段引入标题栏字段级多模态增强。
  \item 输出最终检索结果，并支持 Grad-CAM 可视化和误检分析。
\end{enumerate}

从层次关系看，图纸输入层负责将用户上传图纸和图库图纸统一转化为系统可处理的 PNG 图像；路线选择层负责根据实验模式决定是否启用清洗、掩码聚合、结构重排和 OCR 辅助；特征提取层负责生成视觉向量、结构描述子与文本字段；向量检索层负责完成大规模候选召回；重排层负责对候选结果进行结构相似度和多模态分数修正；结果分析层则面向用户展示检索结果、统计信息和可解释性结果。通过这种分层方式，系统能够在保持在线检索流程简洁的同时，为论文实验提供较充分的模式切换能力。

\section{数据集组织与预处理设计}

根据任务书和开题报告中“构建一万张以上、二十类以上工程图纸数据集”的要求，本文将清理后的 \texttt{Dataset\_img3} 作为正式实验和系统建库数据。该数据集中的图纸均为 PNG 格式，来源于 DWG 图纸转换结果。由于原始转换过程中存在黑底、黑边和截屏边框等问题，系统在正式建库前首先进行统一背景处理，使图纸尽量保持白底线稿风格，从源头降低无关背景对视觉编码器的影响。

在数据组织上，系统按类别目录管理图纸文件，并在建库阶段写入图像路径、类别标签、文件名和结构描述等元数据。这样的组织方式有两点好处：一方面便于前端进行类别浏览和图库统计；另一方面便于离线评测脚本根据类别标签构造查询集合并计算 Recall、mAP、nDCG、FT、ST 和 ANMRR 等指标。对于后续扩展任务，系统也可以在保持目录结构不变的情况下继续加入新的图纸类别。

\section{双主线模式管理设计}

为了避免不同实验模式之间逻辑混杂，系统将检索模式抽象为若干可组合开关，包括是否启用 YOLO 清洗、是否启用掩码引导视觉增强、是否启用结构重排以及是否启用 OCR 辅助。论文中的 8 组实验模式均由这些开关组合得到。该设计使双主线实验具有统一实现基础，也便于在前端演示或离线评测时快速切换模式。

具体而言，\texttt{baseline} 仅使用 CLIP 视觉相似度；\texttt{cleaning\_only} 在视觉编码前增加清洗步骤；\texttt{masked\_pooling} 在清洗基础上加入局部视觉增强；\texttt{full\_model} 继续加入结构重排；\texttt{multimodal\_text} 在清洗和结构增强基础上加入 OCR 辅助。无清洗注意力路线则保留原图，不执行 YOLO 清洗，通过注意力式主体增强与结构重排构成 \texttt{attention\_visual} 和 \texttt{attention\_structure}，并在此基础上形成 \texttt{attention\_multimodal}。这种模式管理方式保证了论文实验结论能够直接追溯到系统运行逻辑。

\section{核心模块划分}

系统划分为以下核心模块：

\begin{enumerate}
  \item 图纸预处理模块：负责图纸格式统一、可选的 YOLO 清洗以及掩码生成。
  \item 特征提取模块：负责视觉嵌入、掩码引导聚合、OCR 文本提取和结构描述子生成。
  \item 向量检索模块：负责图纸向量入库、索引构建和候选召回。
  \item 结构重排模块：负责结构相似度计算与候选精排。
  \item 多模态排序模块：负责标题栏字段匹配和候选前排排序修正。
  \item 实验分析模块：负责正式评测、误检案例生成和可解释性可视化。
\end{enumerate}

\section{数据流与服务流}

从系统运行流程来看，输入图纸首先进入 FastAPI 服务，通过统一接口触发路线判断、特征提取与检索过程；随后视觉向量及相关元数据被送入 Qdrant 完成候选召回；最终结果再由业务层执行结构重排，并在启用多模态模式时叠加 OCR 辅助修正后返回。这种“向量召回 + 业务重排”的体系将高性能 ANN 检索与任务相关规则有效解耦，既有利于系统扩展，也有利于双主线实验的统一比较。

在离线实验流程中，系统不直接依赖在线接口返回结果，而是针对每一种模式分别生成图库特征缓存与查询特征缓存，再在统一评测脚本中完成相似度计算和指标统计。这样做可以避免在线参数切换带来的评测不一致问题，确保图库特征与查询特征始终来自同一检索模式。该流程虽然计算量更大，但更适合作为论文中的正式实验依据。

\section{系统部署思路}

从部署方式看，本文系统采用“算法服务 + 向量数据库 + 前端页面”的组织方式。其核心考虑在于，将模型推理、向量检索与页面展示三个层次解耦，可以降低系统维护复杂度，也有利于后续进行单独扩展。其中，FastAPI 负责统一对外提供检索与管理接口，Qdrant 负责存储向量及元数据，前端页面负责完成图纸上传、结果展示与统计浏览。该部署思路既满足了毕业设计的演示需求，也为实际应用场景中的服务化改造提供了基础框架。

\section{设计原则}

系统设计遵循以下原则：

\begin{enumerate}
  \item \textbf{视觉主导。} 检索结果首先由视觉向量召回保证全局可用性。
  \item \textbf{双路线兼容。} 系统既支持显式清洗路线，也支持无清洗注意力式增强路线，保证方法对比具有统一实现基础。
  \item \textbf{结构增强。} 结构先验作为工程图纸最重要的补充信息，用于提升 Top-K 排序质量。
  \item \textbf{多模态微调。} 标题栏字段仅作用于候选前排重排，避免 OCR 噪声对整体候选集合进行大范围改写。
  \item \textbf{工程可复现。} 所有实验均通过可运行脚本和服务接口完成，便于复现实验结果。
\end{enumerate}

\section{本章小结}

本章从系统角度给出了工程图纸检索平台的总体设计方案，明确了系统目标、整体架构、核心模块、数据流与设计原则。与传统单一路径系统不同，本文系统在架构层面即支持 YOLO 清洗路线与无清洗注意力路线两种模式，并通过统一的向量召回、结构重排和多模态增强框架完成结果输出。这种设计不仅保证了系统运行的完整性，也为后续方法对比实验提供了统一、可复现的实现基础。
